\documentclass[11pt]{article}
\usepackage[T1]{fontenc}
\usepackage{lmodern}
\usepackage[a4paper,margin=25mm]{geometry}
\usepackage{amsmath,amssymb,amsthm}
\usepackage[hypertexnames=false,hidelinks]{hyperref}
\newcommand{\V}{\mathsf V}
\newcommand{\E}{\mathsf E}
\newcommand{\PE}{\Phi_{\E}}
\newcommand{\PV}{\Phi_{\V}}
\newcommand{\flow}{\mathsf{flow}_{+}}
\newcommand{\loc}[1]{\par\smallskip\noindent{\small\textit{Source: #1}}\par\smallskip}
\setlength{\emergencystretch}{2em}
\begin{document}
\begin{center}
{\large Probability on Trees and Networks\\[2pt]Exercise 6.34 and its Author Comment}\\[4pt]
Russell Lyons\quad Yuval Peres\\[2pt]
Version of 19 August 2026---corrected and updated paperback edition
\end{center}

% BEGIN CONTEXT 6.34
\noindent\textbf{Source and scope (editorial).}
This private reader reproduces selected English author text from the fixed 716-page paperback PDF (21,010,205 bytes). Exercise 6.34 is on physical p.\ 246 (printed p.\ 229); its semantic target is the complete unweighted statement of Theorem 6.4 on physical p.\ 199 (printed p.\ 182). Its complete Comment is on physical p.\ 633 (printed p.\ 616). The supporting Theorem 6.3 and its complete printed vertex-capacity global-flow construction are reprinted below as uncounted context. The body theorem and exercise are one target; Exercise 6.35's refinement is excluded. General $D$ in the supporting duality is not added to the primary unweighted theorem.

The book authors are Russell Lyons and Yuval Peres. The Comment identifies this as the original BLPS (1999b) proof; the book is the source transcribed here, and no separate historical paper has been acquired. The original book and the separate errata dated 19 August 2026 remain distinct materials. The errata entry for printed p.\ 616 says to replace both $v$'s by $x$, and to replace the final ``for some $x$'' by ``for some $x\in\partial_{\V}K$''. The fixed book already prints those readings; they are retained, without a new silent correction. No source mathematics is expanded, repaired, or independently certified.

\medskip\noindent\textbf{Oriented-edge and transitivity notation.}
\loc{Section 1.1, physical p.\ 19 (printed p.\ 2); Section 2.4, physical p.\ 49 (printed p.\ 32); transitivity definition, physical p.\ 196 (printed p.\ 179).}
The endpoints $e^-$ and $e^+$ are the tail and head of $e$; $-e$ reverses orientation. An antisymmetric function satisfies $\theta(-e)=-\theta(e)$. The net flow out of $x$ is denoted
\[
(d^*\theta)(x):=\sum_{e^-=x}\theta(e).
\]
The book calls a network transitive when, for every pair of vertices, a graph automorphism takes one to the other and preserves the edge and vertex weights. In the unweighted graph target below, transitivity is the corresponding graph-automorphism condition. No square-summability restriction is added to the antisymmetric functions in Theorem 6.3 or the Comment.

\noindent\textit{Standing chapter setting (physical p.\ 194, printed p.\ 177).} Section 6.1 begins with a connected graph $G$. Its networks have positive edge weights $c(e)$ and vertex weights $D(x)$; $G$ is assumed locally finite or, more generally, $\sum_{e^-=x}c(e)<\infty$ for each $x$, as in (6.1). The primary below remains the original unweighted graph statement.

\subsection*{The outer vertex boundary}
\loc{Physical p.\ 197 (printed p.\ 180); the notation $|K|_D$ is defined on physical p.\ 194 (printed p.\ 177).}
Given positive numbers $D(x)$ on the vertices of a graph $G$, define the (outer) vertex boundary $\partial_{\V}K$
\[
\partial_{\V}K:=\{x\notin K;\ \exists y\in K\quad y\sim x\}
\]
and
\[
\PV(G):=\PV(G,D):=
\inf\left\{\frac{|\partial_{\V}K|_D}{|K|_D};\ \varnothing\ne K\subset\V\text{ is finite}\right\},
\qquad |K|_D:=\sum_{x\in K}D(x),
\]
called the vertex-expansion constant or vertex-isoperimetric constant of $G$. The most common choice is $D=1$. We call $G$ vertex amenable if its vertex-expansion constant is $0$.

\subsection*{The complete supporting vertex-capacity construction}
\loc{Theorem 6.3 and its proof, physical pp.\ 198--199 (printed pp.\ 181--182).}
The analogue for vertex amenability of Theorem 6.1 is due to Benjamini and Schramm (1997). It involves the amount flowing along edges into vertices,
\[
\flow(\theta,x):=\sum_{e^+=x}\bigl(\theta(e)\vee0\bigr).
\]
\medskip\noindent\textbf{Theorem 6.3 (Duality for Vertex Expansion).}
\textit{For any graph $G$ with vertex weights $D$, we have}
\[
\begin{split}
\PV(G,D)=\max\bigl\{\alpha\geq0;\ &\exists\theta\ \forall x\ \flow(\theta,x)\leq D(x)\text{ and }\\
&d^*\theta(x)=\alpha D(x)\bigr\},
\end{split}
\]
\textit{where $\theta$ runs over the antisymmetric functions on $\E$.}
\begin{proof}
This time we use the max-flow min-cut theorem in the version of Exercise 3.14, where capacity constraints are imposed on vertices as well as on edges. Given $\alpha\geq0$ and any finite, nonempty $K\subset\V$, define the network $G(K)$ with vertices $K\cup\partial_{\V}K$ and two extra vertices, $a$ and $z$. The edges of $G(K)$ are those of $G$ that have both endpoints in $K\cup\partial_{\V}K$, an edge between $a$ and each point in $K$, and an edge between $z$ and each point in $\partial_{\V}K$. Give all edges incident to $a$ capacity $c(a,x):=\alpha D(x)$. Let the capacity of the vertices in $K$ be $c(x):=(\alpha+1)D(x)$, and the capacity of the vertices in $\partial_{\V}K$ be $D(x)$. The remaining edges and vertices are given infinite capacity.

In view of the cutset consisting of all vertices in $\partial_{\V}K$, it is clear that any admissible flow from $a$ to $z$ in $G(K)$ has strength at most $|\partial_{\V}K|_D$. Now a function $\theta$ satisfying $\forall x\ \flow(\theta,x)\leq D(x)$ and $d^*\theta(x)=\alpha D(x)$ induces an admissible flow on $G(K)$ from $a$ to $z$ of strength $\alpha|K|_D$, whence $\alpha\leq|\partial_{\V}K|_D/|K|_D$. Since this holds for all $K$, we get $\alpha\leq\PV(G)$.

In the other direction, if $\alpha\leq\PV(G)$, then for all nonempty $K$, we claim that there is an admissible flow from $a$ to $z$ in $G(K)$ of strength $\alpha|K|_D$. Consider any cutset $\Pi$ of edges and vertices separating $a$ and $z$. We will show that $|\Pi|_c\geq\alpha|K|_D$, whence our claim will follow from the max-flow min-cut theorem. If $\Pi$ contains an edge of infinite capacity, then this is trivial, so assume it does not. Let $K'$ be the vertices in $K$ that $\Pi$ separates from $z$. Then $\partial_{\V}K'\subset\Pi$ and $[a,x]\in\Pi$ for all $x\in K\setminus K'$. Therefore,
\[
\begin{aligned}
\sum_{e\in\Pi\cap\E}c(e)+\sum_{x\in\Pi\cap\V}c(x)
&\geq\alpha|K\setminus K'|_D+|\partial_{\V}K'|_D\\
&\geq\alpha|K\setminus K'|_D+\PV(G)|K'|_D\geq\alpha|K|_D.
\end{aligned}
\]
Thus, our claim follows. Note that the flow along $[a,x]$ is $\alpha D(x)$ for every $x\in K$.

Now let $K_n$ be finite, nonempty sets increasing to $\V$, and let $\theta_n$ be the corresponding admissible flows on $G(K_n)$. There is a subsequence $\langle n_i\rangle$ such that for all edges $e\in\E$, the limit $\theta(e):=\theta_{n_i}(e)$ exists; clearly $\forall x\ \flow(\theta,x)\leq D(x)$ and $d^*\theta(x)=\alpha D(x)$.
\end{proof}
\noindent\textit{Source input boundary (editorial).} The vertex-capacitated max-flow min-cut theorem is the named ordinary input from Exercise 3.14. Its general proof is not supplied as a new argument here. The full application, limiting construction, and the Comment's actual strictness argument are retained.
% END CONTEXT 6.34

\clearpage
% BEGIN PRIMARY STATEMENT 6.34
\section*{Exercise 6.34 and its full target}
\loc{Exercise 6.34: physical p.\ 246 (printed p.\ 229).}
\noindent\textbf{Exercise 6.34.} Use Theorem 6.3 and its proof to give another proof of Theorem 6.4.

\loc{Theorem 6.4: physical p.\ 199 (printed p.\ 182).}
\noindent\textbf{Theorem 6.4.}
\textit{If $G$ is an infinite, transitive graph, then for all finite, nonempty $K\subset\V$, we have $|\partial_{\V}K|/|K|>\PV(G)$.}
% END PRIMARY STATEMENT 6.34

% BEGIN PRIMARY PROOF 6.34
\section*{Complete author Comment on Exercise 6.34}
\loc{Physical p.\ 633 (printed p.\ 616).}
(This was the original proof of Theorem 6.4 in BLPS (1999b).) Let $\theta$ be an antisymmetric function on $\E$ with $\flow(\theta,x)\leq1$ and $d^*\theta(x)=\PV(G)$ for all vertices $x$. Let $K\subset\V$ be finite and nonempty. Suppose for a contradiction that $|\partial_{\V}K|/|K|=\PV(G)$. The proof of Theorem 6.3 shows that for all $x\in\partial_{\V}K$, we have $\flow(\theta,x)=1$ and $\flow(-\theta,x)=\PV(G)+1$; in particular, there is some $e$ leading to $x$ with $\theta(e)\geq1/d$ and some $e$ leading away from $x$ with $\theta(e)\geq(\PV(G)+1)/d\geq1/d$, where $d$ is the degree in $G$. Since $G$ is transitive, the same is true for all $x\in\V$. Therefore, we may find either a cycle or a bi-infinite path with all edges $e$ having the property that $\theta(e)\geq1/d$. We may then subtract $1/d$ from these edges, yielding another function $\theta'$ that satisfies $\flow(\theta',x)\leq1$ and $d^*\theta'(x)=\PV(G)$ for all vertices $x$. But then $\flow(\theta',x)<1$ for some $x\in\partial_{\V}K$, a contradiction.
% END PRIMARY PROOF 6.34

\medskip\noindent\textbf{Cached bibliographic attribution (not an additional source proof).}
The book bibliography, physical p.\ 667 (printed p.\ 650), lists: Benjamini, I., Lyons, R., Peres, Y., and Schramm, O.\ (1999b), \textit{Group-invariant percolation on graphs}, \textit{Geom.\ Funct.\ Anal.}\ \textbf{9}(1), 29--66, MR: 99m:60149.
\end{document}
